\begin{titelseite}
\quellseite{7}
\werktitel{CATECHETICS}

\werkuntertitel{Theory and Practise of Religious Instruction}
\end{titelseite}
\begin{titelseite}
\quellseite{9}
\werktitel{CATECHETICS}

\werkuntertitel{or Theory and Practise of Religious Instruction}\\[1.5em]

By M. REU, D.D., Lit. D.\\
Professor of Theology at Wartburg Seminary\\
Dubuque, Iowa\\[1.5em]
SECOND, REVISED EDITION

CHICAGO\\
WARTBURG PUBLISHING HOUSE\\
1927
\end{titelseite}

\quellseite{11}
\chapter{Preface to the First Edition}
THIS textbook on Catechetics was originally published in German, the first edition appearing in 1915. Since it was introduced by three Lutheran seminaries a second edition soon proved necessary and was published in the spring of 1918. Ever since its first publication the desire was often expressed for an English edition. Thanks to the efforts of Rev. C. B. Gohdes, Litt. D., and Professor of History at the Capital University, Columbus, Ohio, this desire is now fulfilled. Dr. Gohdes kindly carried out the very difficult task of translation. In rendering the chapters 21 and 22 he was kindly assisted by Prof. Theo. Mees, Ph.D., D.D., of Capital University, whose conversance with the psychological terminology of the Herbart Society was placed at his disposal, while Rev. Harry Melcher of Columbus, Ohio, furnished the translation of the Practical Examples. For all such valuable assistance, especially for the unselfish, assiduous, and successful labors of Dr. Gohdes, the author takes this opportunity publicly to express his heartfelt thanks.

Although the esteemed translator, under constant collaboration with Rev. S. Salzmann, of Dubuque, Iowa, and myself, gave much time and diligent thought toward the production of this English edition, here and there an error may have crept in. May I not ask the readers to make a note of these, and kindly to bring them to my attention in order that they might be corrected in case a second edition should be necessary? I should likewise be grateful for any other suggestions.

This book, however, is not merely a translation of the German edition. For, in its present form, it was not only divested of everything which had special reference to conditions
\quellseite{12}
in the German speaking Lutheran congregations of our country, but was also recast and amplified, as this seemed necessary for a more general use in our Church. While already in the German edition the psychological, pedagogical, and catechetical literature of our country was drawn upon, especially for chapter 22, this has now been registered and turned to account throughout the book. Chapters 30 and 31 dealing with the educational agencies and the distribution of material has been completely recast and amplified. Chapter 20 appears for the first time. It seemed necessary, in order that the student and pastor might, at least briefly, be informed on the work of religious instruction carried on by the churches round about us. One addition which I had planned, could not be embodied in this print. I had intended to write a sketch of the catechetical work of the Scandinavian Church from the Reformation to the present time, but being somewhat overworked, I had to desist from this plan in the last moment.

To some it might seem as if this Catechetics were too bulky to serve as a textbook; it was, however, purposely written in this volume, for it should be of service also to the pastor and catechist, who have stood in the ministry for years. By omitting the literature and the practical examples about one hundred pages could have been eliminated. But perhaps just the teacher of Catechetics and all who intend later to continue the study of catechetical problems, will thank me for having recorded it with such care. Hardly will they find one so relatively complete. Most of the books mentioned are found in my library, the remainder was at least examined, only a few have been cited on the authority of others. If a large proportion of the literature given is written in German, the reason for this is mainly, because nowhere has such diligent, practical as well as scientific, work been done in the catechetical field as in the Evangelical Church of Germany. However much conditions may change, it ought to remain a rule with us that
\quellseite{13}
the Lutheran pastor or, at least, the Lutheran professor should be able to read and to understand the language in which Luther wrote. The practical examples could, in the interest of the catechist, especially of the student and beginner, not well be omitted. While some of the examples given, in my judgment, can by no means serve as models worthy of imitation, they have been added, because, in the form given, they have proved valuable for my work in the Seminary. Here, too, criticism of the material offered serves as an aid to find and appreciate what is correct. Another reason why the bulkiness of this volume should be no hindrance to its usefulness as a textbook is this: in spite of its unity the several parts are complete in themselves and can, therefore, also be taken up and studied separately. If time will not allow for a careful study of the historical part, this can be covered by a few summarizing lectures. He who has reasons to believe that his class is sufficiently grounded in the elements of psychology may restrict himself to a hasty review of the second part in which stress is laid upon the pedagogical deductions. My experience, however, is that such knowledge of psychology is rather rare among our students of theology. Where this is the case much of the necessary basis for all catechetical work is wanting.

I can truthfully say that this textbook has grown out of scientific as well as practical study of catechetical problems extending over many years. Especially what is said concerning the various educational agencies and the distribution of material has been tested as to its practicableness either by myself or by some of my former pupils who perform all their catechetical work in English.

Would, that now, when the transition into the American language has made religious instruction much easier in many respects, all would earnestly strive to “redeem the time” and to lay the foundations for a lasting improvement of ways and means of such instruction. The whole future of our Church
\quellseite{14}
is at stake. In the same measure in which we succeed in solving the question of the religious instruction of our youth, can we become what we should, a salt for our country. The following resolution was passed by the International Sunday School Convention assembled this summer at Buffalo: “The prize of our religious liberty is the sum required for the building of a system of Church schools which will parallel our system of public schools and be equally efficient. We do not have in this country a system of public education; we have only a system of public schools; but this system of schools does not work with the whole child. It is but half an educational arch. We must complete the arch by building a system of Church schools closely co-ordinated with the public schools. These two systems of schools—one supported by the State, with secular leadership, the other supported by the Church, with religious leadership—will form the only system of education that a country can have, in which the Church and the State are a part. The building of this system of Church schools is the task now pressing for completion.” At the same time the following report comes from New York: “The Commissioner of Education for the State of New York declares that there must be some definite plan of religious education for the children of the State, this to be formulated through the co-operation of the schools and the Churches. He proposes three methods: (1) The preparation, for use in the schools, of a book of selections from the Bible by an interdenominational commission appointed by the legislature; (2) the formulation of a plan for co-operation between the school and the various denominations, that every child may be provided with religious instruction; (3) the granting of regents’ credits for serious work in Bible study outside of the schools.”

I still maintain what has been said in chapters 20, 30 and 31 concerning the establishment of parish schools, but in connection with these latest utterances I would like once more to
\quellseite{15}
emphasize the following: Where parish schools can not be established, there let us unite and strive with all seriousness that a portion of the regular school periods of our public schools be assigned for religious instruction given by the Church, and let us train teachers for that purpose. To bring this matter before our state legislatures and boards of education would be a very important task for the newly formed National Lutheran Council. In the meantime let no one idly wait for developments which the future might bring, but let him turn to account the educational agencies mentioned on page 402, or let him choose other ways and means—if only our youth receives that religious instruction which the Church and every pastor is in holy duty bound to give.

\begin{unterschrift}
Dubuque, Iowa, September, 1918.\\
M. REU.
\end{unterschrift}

\quellseite{17}
\chapter{Preface to the Revised Edition}
This new edition varies but little from the first as far as contents are concerned. Wherever it was necessary, I have made changes or additions to bring the volume up to date. The most important recent publications in the field of Christian Education have been listed in their proper place; and upon repeated requests I have marked with an asterisk those works which are of special value. The translation, however, has been thoroughly revised, and many pages have even been entirely rewritten. This work was done by my colleague, Professor Julius Bodensieck, who at the present time teaches Religious Education in our theological seminary. The indexes were, likewise furnished by him, and he has seen the volume through the press. I would also at this place express my gratitude to him for the thoroughness and devotion with which he performed his arduous task.

\begin{unterschrift}
Dubuque, Iowa, January 10, 1927.\\
M. REU.
\end{unterschrift}

\quellseite{21}
RE—Realenzyklopaedie fuer protestantische Theologie und Kirche in dritter Auflage herausgegeben von A. Hauck.

Tr.—Triglot Concordia. The Symbolical Books of the Ev. Lutheran Church, St. Louis, 1921.

Page 319, line 17, read First Chief Part.
